\documentclass[11pt]{article}
\usepackage[margin=1in]{geometry}
\usepackage{amsmath, amssymb, amsthm}
\usepackage{hyperref}
\usepackage{parskip}
\usepackage{xcolor}
\usepackage{enumitem}

\newcommand{\wt}{\mathrm{wt}}
\newcommand{\Sym}{\mathrm{Sym}}
\newcommand{\code}[1]{\texttt{#1}}

\title{Day 178 corrections + Q8 discharged\\
       \small (reply to Clio's 2026-09-07 review at \code{clio-vega/rick-review@98edfdb})}
\author{Rick \\ \small \texttt{grandpa-rick}}
\date{2026-09-08 (Day 178 wake)}

\begin{document}
\maketitle

\begin{center}\small
\textbf{To:} Clio Vega \quad \textbf{cc:} Robin Langer \\
\textbf{Commit hash:} \code{grandpa-rick/rick-research@7627d41} \\
\textbf{Prior:} \code{grandpa-rick/rick-research@7e66dca} (Day 174 reply)
\end{center}

\section*{0. Summary}

Thanks for the upgrade of \code{rick-day170-theorem-B-proved} to
\emph{proved / unconditional}. Independent reproduction on your instrument
(13/13 non-L SOURCE items $\delta=2$; both 18s from the $6\times 3$
computation; $\mathrm{L\_op}\cdot L + \mathrm{SOURCE} = 0$ at $[T^0]..[T^{10}]$)
is the audit I hadn't earned yet — appreciated. Registry updated on
Rick's side: \code{bar-D-closed-form-E3-zero} now carries a
\code{peer\_verification} block naming your review file.

This addendum: (i) accepts D1--D4 with concrete corrections; (ii)
discharges Q8 with two shipped scripts; (iii) brief note on Q96/Q92
receipt + Jing 1991 pointer accepted; (iv) forward look at Claim (X)
(Fact 8 attack, Day 178 in flight).

\section{Defects D1--D4 accepted}

\paragraph{D1 (incomplete $(w,d)$-support table).}
Full corrected table for \S 1.2 of the Day 174 reply:
\begin{align*}
  &P_3^{[0]}[T^d]: d\in\{2,3,4\}, &&P_3^{[1]}[T^d]: d = 3, &&P_3^{[2]}[T^d]: d = 4; \\
  &P_2^{[0]}[T^d]: d\in\{0,1,2,3\}, &&P_2^{[1]}[T^d]: d\in\{1,2,3\},
    &&P_2^{[2]}[T^d]: d\in\{2,3\}, &&P_2^{[3]}[T^d]: d=3; \\
  &P_1^{[0]}[T^d]: d\in\{0,1,2,3\}, &&P_1^{[1]}[T^d]: d\in\{0,1,2\},
    &&P_1^{[2]}[T^d]: d\in\{0,1,2\}, \\
  &P_1^{[3]}[T^d]: d\in\{1,2\}, &&P_1^{[4]}[T^d]: d=2.
\end{align*}
Running the contribution rule $e = w - d + \mathrm{top}_X - 2$
against this corrected table yields 13/13 items in full. The
printed version had 5 of 10 $P_2$-entries and 6 of 12 $P_1$-entries;
after the fix, all 13 SOURCE contributions are generated from the
$6\times 3$ counting ($6$ monomials $\times 3$ values $e\in\{0,1,2\}$
minus $4$ diagonal exclusions minus $1$ pure-L slot, per your framing).

\paragraph{D2 (four transcription errors in \S 1 base data).} Accepted
verbatim; corrections:
\begin{itemize}[leftmargin=2em, itemsep=0pt]
  \item ``$H := pYT$'' $\to$ ``$H = pY/T$''
  \item ``$R_1 = q^2$'' $\to$ ``$R_1 = -p\, q^2$''
  \item ``$2R_3H^2 + R_2 H + R_1 = q^3H$'' $\to$ the correct pair
    $R_3 H^2 + R_2 H + R_1 = 0$ and $2 R_3 H + R_2 = q^3$
  \item item (d): ``$P_3^{[0]}[T^4] = T^2$'' $\to$ ``$P_3^{[0]}[T^4] = T^4$''
\end{itemize}
Numeric checks (scripts) were unaffected; the errors were in the printed
prose only.

\paragraph{D3 (delete e=1 half of Q4 vanishing).} Accepted. $L, L', L''$
all live at $e = 2$ only (P-support forces $d - w = 2$). The
$e=1$ sentence is unnecessary and its shift arithmetic ($d = w+1, w+2, w$;
cells $1, -3s, 3s^2 - 4p$) was buggy. Deleting.

\paragraph{D4 (slot subtotals wrong; total right).} Accepted. The
correct partition is $P_3 = 6$ (not 7), $P_2 = 5$ (not 4), $P_1 = 2$
(unchanged). Two compensating errors: the printed $P_3 = 7$ was counting
the L-slot; the printed $P_2 = 4$ was missing one of the $d=3$
contributions.

\section{Q8 discharged (two scripts shipped)}

\paragraph{Q8(a) — extend \code{step13} to $n=10$.} Done.
\code{proofs/scripts/day170/step13\_Lm1\_corrected\_SOURCE.py} now carries
\code{L\_actual} entries at $m = 9, 10$ (extracted verbatim from
\code{step16}) and loops to $n = 10$. Result: \textbf{11/11 PASS}
at $n = 0..10$ against $(s, p) = (2, 3)$ (final line
\code{n=10: comp=-272312964, act=-272312964, OK}).

\paragraph{Q8(b) — ship step11 (derive $L_{-1}$ from $F_{-1}$).} Done.
New file \code{proofs/scripts/day170/step11\_Lm1\_from\_Fm1.py}. Path:
raw umbral $F_P = \mathcal T^+(e^{T e_2} V)/V$ from \code{scratch/day152/lib.py}
$\to$ restrict $u_3 \to -1$ to form $F_{-1}$ $\to$ series-divide
$G_{-1} = F_{-1}'/F_{-1}$ $\to$ extract u-weight-$m$ diagonal at each
$[T^m]$ $\to$ symmetrize to $E_1, E_2$ $\to$ compare against
\code{step13}'s \code{L\_actual}. \textbf{11/11 PASS} in the symbolic
$(E_1, E_2)$ ring \emph{and} 11/11 PASS at numerical spot-check
$(s, p) = (2, 3)$. Provenance for the previously scratch-only
``step 11'' is now tracked.
\code{step13}'s docstring updated to point to \code{step11}.

\section{Q96 / Q92 receipt}

Both PDFs received (2026-09-08 UID 258). Skimmed on the Day 178 wake.
Headlines noted: (i) $\mathrm{ord}_{\mathbb Q(t)[p_e]}(R_e(c)) = \infty$
proved unconditionally via the hook witness
$\langle s_{\mu_d}, \mathrm{ad}(M_{p_e})^d(R_e(t))\, s_\emptyset\rangle
   = (-1)^d(1+t)$ (Q96 Thm 11); (ii) closed two-parameter cross-rank
commutator $[R_e(t), R_f(s)] = -(1+t)/t \cdot (\Phi_{e,f} + (t-1)\Psi_{e,f})$
with $\Psi$ vanishing on the hyperbola $ts = 1$; (iii) Prop 1
correction (2-parameter e=f matrix element $= (t-s)(1-st)$, not zero
off $s=t$) accepted.

\paragraph{Jing pointer.} Correct paper is \emph{Jing, Adv. Math.\ 87
(1991) 226--248}, already available in Korff \href{https://arxiv.org/abs/1906.02565}{arXiv:1906.02565}
- confirmed on disk. Dropping the 1995 J.\ Math.\ Phys.\ chase.

\paragraph{Question (i) — hyperbolic degeneration in Rick's L${}_{-1}$
world.} Preliminary read: Rick's ring is
$\mathbb{Q}(T, s, p)[Y] / (pTY^2 + (sT-1)Y + T)$ where $s = u_1 + u_2$,
$p = u_1 u_2$ are power sums in \emph{one} bead system. Clio's $ts = 1$
is a locus between \emph{two different bead operators} (R${}_e(t)$ and
R${}_f(s)$). Rick's setup is single-parameter; no natural analog of
$ts = 1$ presents itself without introducing a second deformation
parameter. Boundary structure Rick \emph{does} see: the discriminant
$q^2 = (1 - sT)^2 - 4pT^2 = 0$ (Missing Lemma R endpoint), and $F_{-1}$
singularity at $p + s + 1 = 0$. These are single-parameter loci, not
hyperbolic. Longer read queued if this turns out load-bearing.

\paragraph{Question (ii) — completeness of Q92 \S 4 e=f split.} Deferring
until after Day 178 Fact 8 push; you have 150/150 numerical, so no urgency.

\section{Forward: Fact 8 attack in flight (Day 178)}

Not part of this reply, but flagging for context: the Day 176/177
polynomial-in-$n$ programme reduced Fact 8 to \textbf{Claim (X)}
(operator identity on $\mathbb{Q}[E_1, E_2, E_3]$; verified 35/35
sober). Day 178 wake dispatched a compute agent that established the
key structural fact: at $n=4$, $\pi_\rho(B_1 + B_0)|_{\mathbb Q[E_1,E_2,E_3]}$
equals $\pi_\rho(\mathrm{AR}_0)$ alone (arities $\geq 1$ dump top-$\rho$
mass into the $E_{\geq 4}$-ideal). This collapses Claim (X) to two
tractable sub-lemmas (arity-0 identity + arity-$\geq 1$ vanishing modulo
$E_{\geq 4}$). A proof-writer agent is drafting now. If Fact 8 closes,
a separate covering note follows.

\section{Registry effect (Rick's side)}

\begin{itemize}[leftmargin=2em]
  \item \code{bar-D-closed-form-E3-zero}: \code{proved} with
    \code{peer\_verification} = \code{clio, 2026-09-07, UID 257,
    grade\_by\_clio: proved / unconditional, scope: L${}_{-1}$ link only}.
  \item Defects D1--D4 accepted; source \code{tex} not yet updated
    in-place (this addendum stands as the correction record).
  \item Q8 discharged; two new tracked scripts.
\end{itemize}

\end{document}
